\section{INTRODUÇÃO}

A expressão Dívida Técnica (DT) foi originalmente cunhada para justificar a necessidade de refatoração de código em projetos de software~\cite{11,14}. Por meio dessa metáfora, ilustrava-se que um código funcional, porém mal estruturado, assemelha-se a uma dívida financeira: com o passar do tempo, sua quitação torna-se obrigatória para evitar a degradação e o comprometimento do desenvolvimento~\cite{14}. Com a evolução das pesquisas na área, o escopo do conceito foi ampliado para englobar novos passivos, tais como conjuntos de testes incompletos, documentação ausente, incompleta ou defasada, além de escolhas arquiteturais inadequadas que prejudicam a qualidade interna do produto~\cite{11,12,13,15}.

Nesse sentido, a contração de uma DT em um projeto ocorre quando algum processo ou artefato essencial do ciclo de desenvolvimento é negligenciado. As causas para essa omissão são diversas, englobando desde o pragmatismo e a priorização de outros objetivos de negócio até a desatenção, falhas nos processos corporativos e fatores comportamentais dos indivíduos envolvidos~\cite{12}. Consequentemente, o acúmulo dessa dívida acarreta impactos negativos no moral e na produtividade da equipe, além de comprometer a qualidade do software e elevar os riscos do projeto~\cite{12}. Contudo, a literatura da área apresenta diversas estratégias para o gerenciamento das DTs, viabilizando a sua mitigação por meio da adoção de práticas consolidadas ao longo do ciclo de vida do software. Para um aprofundamento na temática, destacam-se pesquisas como o mapeamento sistemático conduzido por Li, Avgeriou e Liang~\cite{13}, voltado para a identificação e o manejo das DTs, e o estudo terciário de Nicolli, de Mendonça Neto e Spínola~\cite{15}, que analisa as tendências acadêmicas referentes ao gerenciamento desse fenômeno.

Dentre as diversas práticas recomendadas para a mitigação e a prevenção da DT, a elaboração adequada da documentação de software destaca-se como um elemento fundamental. De modo geral, documentar um projeto consiste no registro formal de suas informações, com o intuito de detalhar tanto o processo de engenharia quanto o funcionamento do produto para os seus usuários~\cite{4}. Contudo, dada a vasta gama de artefatos referentes a essa área, o presente estudo concentra-se especificamente na documentação de sistema. Conforme a definição de Sommerville~\cite[tradução nossa]{1}, essa categoria compreende ``a descrição do produto do ponto de vista dos engenheiros desenvolvendo e mantendo o sistema'', a qual abrange ``todos os documentos descrevendo o sistema desde a especificação de requisitos até o plano final de testes de aceitação".

Um exemplo representativo da documentação de sistema é o documento de requisitos, concebido por meio da aplicação de técnicas da engenharia de requisitos. Essa disciplina compreende o processo de elicitação e registro de informações~\cite{20}, com o propósito de especificar tanto os serviços a serem disponibilizados pelo software --- definidos como requisitos funcionais --- quanto as restrições impostas à sua operação --- classificadas como requisitos não funcionais~\cite{8}. Adicionalmente, esse artefato costuma registrar as regras de negócio (RNs), que consistem em diretrizes estabelecidas para determinar ou influenciar a execução das atividades da organização~\cite{23}. Desse modo, a consolidação dessas três frentes de especificação resulta em um instrumento essencial, utilizado para orientar desenvolvedores e engenheiros de software quanto ao comportamento esperado do sistema, às suas limitações operacionais e às lógicas do domínio no qual está inserido.

A adoção das práticas de documentação de requisitos proporciona diversos benefícios ao projeto de software. Dentre eles, destacam-se: a melhoria na compreensão do sistema em desenvolvimento; o suporte contínuo às fases de implementação e manutenção; a redução do tempo despendido na execução de tarefas operacionais; e o embasamento para a tomada de decisões relativas ao produto~\cite{3}. Contudo, o elevado esforço inerente à sua manutenção --- seja ele expresso em tempo ou recursos financeiros --- frequentemente motiva a negligência dessas técnicas~\cite{3, 2}. Esse cenário induz a equipe de desenvolvimento a priorizar ganhos imediatos, mesmo ciente de que tal escolha implicará prejuízos futuros. Incorre-se, assim, na contração de uma DT de documentação.

Devido ao efeito cascata inerente à DT, a negligência documental propaga-se rapidamente para outras fases do ciclo de desenvolvimento, resultando na contração de novas dívidas. No contexto específico dos testes de software, a finalidade dessa prática é, segundo Sommerville~\cite[p. 227, tradução nossa]{8}, ``mostrar que um sistema faz o que ele deveria fazer e revelar os seus defeitos antes de ele começar a ser utilizado''. Contudo, ao se contrair uma DT na documentação de requisitos, a simples definição do que o sistema ``deveria fazer'' torna-se uma tarefa complexa e ambígua. Consequentemente, a elaboração de casos de teste é severamente prejudicada pela falta de diretrizes exatas sobre o comportamento esperado do produto. Esse cenário culmina na introdução de novas DTs diretamente na etapa de testes, o que inviabiliza a garantia da corretude e da qualidade geral do sistema.

Dentre os mecanismos explorados na literatura para verificar a incidência de DTs na etapa de testes de um projeto, destaca-se a análise de cobertura. Essa prática consiste na mensuração do grau em que um item de software --- o qual pode abranger desde um único requisito até um componente inteiro do sistema --- é verificado ou validado pelos casos de teste~\cite{21}. Desse modo, a aplicação dessa técnica permite avaliar, por exemplo, a proporção de requisitos do sistema que está, de fato, amparada pelas especificações de teste. Contudo, diante da existência prévia de uma DT na documentação de requisitos, torna-se imperativa a mitigação desse passivo antes que qualquer avaliação de cobertura possa ser conduzida de maneira eficaz.

\subsubsection{Objetivos}

O presente trabalho tem como objetivo relatar a experiência de análise e reestruturação da documentação de requisitos e de testes em um projeto real de desenvolvimento, o Sistema de Fomento ao Microempreendedor (SFM), diagnosticando os impactos da DT na etapa de verificação e validação do \textit{software} e nas rotinas da equipe.

\subsubsection{Objetivos específicos}

Para que o objetivo geral seja alcançado, foram traçados os seguintes objetivos específicos:
\begin{enumerate}[topsep=0em]
    \itemsep=0.3em
    \item \textbf{Elaborar} um documento formal de requisitos para o SFM;
    \item \textbf{Mapear} a relação entre os requisitos elicitados e os testes de sistema existentes por meio da construção de uma Matriz de Rastreabilidade;
    \item \textbf{Avaliar} a completude dos casos de teste especificados para o sistema por meio de uma análise de cobertura, evidenciando as lacunas de verificação em ambas as frentes;
    \item \textbf{Investigar} a percepção da equipe técnica quanto ao acompanhamento da evolução do sistema e à relevância da especificação formal de requisitos;
    \item \textbf{Mitigar} parte das DTs contraídas ao longo do ciclo de desenvolvimento.
\end{enumerate}

\subsection{Estrutura do relatório técnico}

Para uma melhor compreensão e fluidez da leitura, este trabalho encontra-se organizado em seis seções principais, além das referências bibliográficas. A Seção 1 introduz o contexto geral da pesquisa, o problema a ser investigado e os objetivos propostos. A Seção 2 aborda os conceitos teóricos necessários para o embasamento do estudo e apresenta uma breve revisão da literatura. A Seção 3, por sua vez, detalha o contexto de desenvolvimento do SFM e descreve a metodologia aplicada para a condução das atividades. A Seção 4 expõe os resultados obtidos e promove uma análise crítica das observações realizadas. A Seção 5 reúne as considerações finais, as contribuições da pesquisa e as sugestões para trabalhos futuros. Por fim, a Seção 6 apresenta a declaração formal sobre o uso de ferramentas de Inteligência Artificial ao longo da elaboração deste documento.